\begin{abstract}
Automatic text-to-query generation has progressed alongside large language models, yet in enterprise settings it still exhibits systematic failures in structural fidelity: references to non-existent tables or columns, invalid joins, and graph patterns inconsistent with the schema. These failures intensify under the hybrid \ac{SQL}/\ac{PGQ} standard (ISO/IEC 9075-16:2023), where generation must operate simultaneously over a relational schema and a property graph defined on top of it.

This thesis proposes a structural reasoning architecture based on a typed intermediate representation, IR-SQL/PGQ, that mediates between the language model's output and the executable query. The architecture articulates four components: a dual-schema grounding stage that narrows the search space; the intermediate representation itself, uniformly expressive over the relational, graph, and hybrid regimes and decidable through local typing rules; a deterministic compiler to \ac{SQL}/\ac{PGQ}; and a cascaded verification loop that operates on the intermediate representation with structured feedback through categorical descriptors.

Empirical validation will be conducted on the Spider and BIRD benchmarks for the relational regime and on a purpose-built corpus for the hybrid \ac{SQL}/\ac{PGQ} regime. Reported metrics will include correctness, hallucination rate, transferability across heterogeneous environments, and aggregated per-query cost. The central hypothesis, subject to empirical verification, is that mediation through a typed intermediate representation with cascaded verification reduces referential and structural hallucinations relative to direct generation, without compromising sustained cost in enterprise scenarios.
\end{abstract}
